\section{Deformation to the normal cone: standard mechanism and the computed special fibre}
\label{sec:rees-specialization-v127}

The passage from an ideal to its extended Rees algebra is the standard
deformation-to-the-normal-cone construction; see
\cite[\S5.1]{FultonIntersection} and
\cite[Tag 062Z]{StacksNormalCone}.  We use it here only as the
canonical mechanism that packages the manuscript-specific graded
algebra computed in Sections~\ref{sec:stratified-nilpotent} and
\ref{sec:boundary-atlas-proofs}.

\begin{proposition}[Intrinsic deformation to the normal cone]
\label{prop:intrinsic-dnc-v127}
Let \(\widehat{\mathcal N}_R\) be the intrinsic nilradical of
\(\widehat D_R\).  On an affine open
\(\Spec A\subset\widehat D_R\), with nilradical \(N\subset A\), put
\[
 \mathscr R_N=A[t,Nt^{-1}]\subset A[t,t^{-1}].
\]
These algebras glue and define a flat family
\[
 \mathfrak D_R\longrightarrow\A^1.
\]
Over \(\mathbb G_m\) it is
\(\widehat D_R\times\mathbb G_m\), while its special fibre is
\[
 \mathfrak D_{R,0}
 \simeq
 \Spec_{\widehat\Delta_R}
 \gr_{\widehat{\mathcal N}_R}
 \mathcal O_{\widehat D_R}
 =
 C_{\widehat\Delta_R/\widehat D_R}.
\]
The construction is functorial for isomorphisms of the abstract
failure scheme.
\end{proposition}

\begin{proof}
The extended Rees algebra is the standard construction.  Multiplication
by \(t\) is injective because \(\mathscr R_N\) is a subring of
\(A[t,t^{-1}]\).  Hence it is torsion-free over the PID \(\C[t]\),
and therefore flat.  Inverting \(t\) gives
\(A[t,t^{-1}]\).  Modulo \(t\), the class of
\(n t^{-j}\), \(n\in N^j\), depends only on
\(n\bmod N^{j+1}\), giving
\[
 \mathscr R_N/(t)\simeq\bigoplus_{j\ge0}N^j/N^{j+1}.
\]
Localization compatibility glues the construction.  Since \(N\) is
the ideal of the reduction \(\widehat\Delta_R\) in
\(\widehat D_R\), the special fibre is the normal cone.  The
nilradical is intrinsic, so the construction is functorial.
\end{proof}

The point of Revision 127 is not Proposition~\ref{prop:intrinsic-dnc-v127}
itself.  The new content is that its special fibre is now bounded and
stratified:
\[
 \gr_{\widehat{\mathcal N}_R}\mathcal O_{\widehat D_R}
 =
 \mathcal O_{\widehat\Delta_R}
 \oplus
 \bigoplus_{j=1}^{5}
 \mathcal L^{\otimes j}\otimes\mathcal O_{W_j},
\]
with the quotient multiplication
\eqref{eq:atlas-multiplication}.  Projection corank two has the
mixed-kernel atlas of
Theorem~\ref{thm:corank-two-mixed-kernel-atlas}; all \(A_H\)-rank
drops have the finite valuation bounds of
Proposition~\ref{prop:rankdrop-coranktwo-depth}; projection corank
three has the exact index-five/index-six dichotomy of
Theorem~\ref{thm:corank-three-depth}; and projection corank four has
exact index five by
Theorem~\ref{thm:corank-four-jacobian}.

Flattening is used only to organize coefficient specializations of
these explicitly presented finite modules.  We use the standard
flattening-stratification formalism in
\cite[Tag 052F]{StacksFlattening}; no novelty is claimed for that
formalism either.

\begin{remark}[Novelty boundary]
The standard ingredients in this section are the extended Rees
algebra, deformation to the normal cone, and flattening
stratification.  The manuscript-specific statements are the Fitting
factorization, the global determinant-power bound \(d^5\in J\), the
nine mixed-kernel normal forms and their graded supports, the
corank-three Schur-type exclusion, the corank-four
Jacobian-covariant calculation, the collision primary laws, and the
polarized reconstruction theorem.  These roles are kept separate in
the statements and in the referee response.
\end{remark}
